\documentclass[11pt,a4paper]{article}
\usepackage[margin=25mm]{geometry}
\usepackage[T1]{fontenc}
\usepackage{lmodern,amsmath,amssymb,amsthm,mathtools,booktabs,array,longtable}
\usepackage{microtype,xcolor,hyperref,fancyhdr}
\hypersetup{colorlinks=true,linkcolor=blue!45!black,urlcolor=blue!45!black,
 citecolor=blue!45!black,pdftitle={Valuation separation and signed certificates for Diophantine consistency},
 pdfauthor={AI-assisted research prepared for Vladimir Reshetnikov}}
\setlength{\headheight}{14pt}
\pagestyle{fancy}
\fancyhf{}\fancyhead[L]{\small Diophantine consistency: two arithmetic reductions}
\fancyhead[R]{\small ProveIt}\fancyfoot[C]{\thepage}
\setlength{\emergencystretch}{2em}
\newtheorem{theorem}{Theorem}[section]
\newtheorem{lemma}[theorem]{Lemma}
\newtheorem{proposition}[theorem]{Proposition}
\newtheorem{corollary}[theorem]{Corollary}
\theoremstyle{remark}\newtheorem{remark}[theorem]{Remark}
\newcommand{\Z}{\mathbb Z}\newcommand{\Q}{\mathbb Q}\newcommand{\N}{\mathbb N}
\newcommand{\Con}{\operatorname{Con}}\newcommand{\ZFC}{\mathrm{ZFC}}
\newcommand{\vp}{v_p}\newcommand{\vt}{v_3}
\newcommand{\Fn}{\mathcal F_{32}}
\newcommand{\Dm}{D_{\mathrm m}}\newcommand{\Df}{D_{\mathrm f}}
\newcommand{\Dt}{D_{\tau}}\newcommand{\Dg}{D_{\mathrm g}}
\newcommand{\Cc}{C_{\mathrm y}}
\newcommand{\Acone}{A_{\mathrm{cone}}}\newcommand{\Bcone}{B_{\mathrm{cone}}}
\newcommand{\primary}{P_{2400}}\newcommand{\companion}{P_{2912}}
\newcommand{\DirectVariables}{1,190}
\newcommand{\DirectMonomials}{26,411}
\newcommand{\DirectHeight}{83,900,628}

\title{\textbf{Valuation separation and signed certificates}\\[3pt]
Explicit Diophantine forms of $\Con(\ZFC)$}
\author{AI-assisted research prepared for Vladimir Reshetnikov\\
\small A reproducible contribution to the ProveIt project}
\date{23 October 2026, Pacific time\\[-1pt]
\small Repository snapshot \texttt{85ae4070019efd66ea3e58e9a67d287528094e8b}}
\begin{document}
\maketitle
\begin{abstract}
We give two explicit improvements to Diophantine representations of syntactic
ZFC refutability. The first is an integer-coefficient polynomial in twelve
unrestricted integer variables, of exact degree $2400$, evaluated by $215$
live binary arithmetic operations. It improves the inspected twelve-integer
construction of degree $4832$ and $248$ operations. A companion has degree
$2912$ and $209$ operations. The main observation is that the existing
interpolated Pell certificate already separates its radicals at the prime
$3$. Constant weights replace growing polynomial weights. Retaining the
auxiliary Pell root as a rational square root, then clearing its denominator
inside a norm, gives a second degree reduction. A classical intersective
quintic can replace the previous monotone carrier.

The second construction is a fully expanded quartic in \DirectVariables{}
unrestricted integers, obtained from the recent synchronized direct
proof-certificate construction. Its coefficients are ordinary decimal
integers of absolute value at most \DirectHeight{}. It uses the calibrated
numerical proof calculus directly, independently of K215 and Jones's
universal polynomial.

For each supplied equation, composing the proved arithmetic transformations
with the identified inherited ordinary calibration proofs establishes,
within ZFC, that absence of solutions is equivalent to $\Con(\ZFC)$.
Neither consistency nor inconsistency is assumed. These are ordinary
mathematical proofs and reproducible finite certificates; no new
proof-assistant derivation of the complete internal ZFC biconditional is
claimed. The contribution includes exact equations, coefficient recipes,
independent checks, source provenance, and further research questions.
\end{abstract}
\clearpage
\tableofcontents
\newpage
\section{Results, comparison, and meaning of the claim}\label{sec:results}

Fix the usual syntactic statement
\[
 \Con(\ZFC)\equiv
 \neg\exists p\in\N\;\operatorname{Proof}_{\ZFC}(p,\bot).
\]
The task is to specify a fixed polynomial $P$ with integer coefficients and
prove
\begin{equation}\label{eq:target}
 \ZFC\vdash
 \left(\neg\exists\mathbf x\in\Z^n\;P(\mathbf x)=0\right)
 \longleftrightarrow \Con(\ZFC).
\end{equation}
Proving this equivalence does not prove either side. In particular, the
construction does not evade the second incompleteness theorem by purporting
to prove its own polynomial has no zero.

\begin{theorem}[Supplied equations]\label{thm:main}
The package specifies polynomials $\primary,\companion\in
\Z[j,C,z,h,r,v,w,T,c,n,t,s]$ with exact degrees $2400$ and $2912$,
respectively. Their delivered straight-line circuits use $215$ and $209$
live binary operations. Each satisfies \eqref{eq:target} with $n=12$,
through the ordinary inherited native/compiler calibration identified in
Section~\ref{sec:dependencies}.

It also supplies a literal quartic $Q_{\mathrm{direct}}$ in
\DirectVariables{} unrestricted integer variables with
\DirectMonomials{} collected nonzero monomials and coefficient height
\DirectHeight{}. This polynomial satisfies \eqref{eq:target} through the
independent numerical-calculus calibration in Section~\ref{sec:direct}.
\end{theorem}

\begin{table}[ht]
\centering\small
\begin{tabular}{@{}llrrr@{}}
\toprule
Construction & Domain & Variables & Degree & Operations\\
\midrule
Prior interpolated norm & $\Z$ & 12 & 4832 & 248\\
Prior shorter evaluator & $\Z$ & 12 & 16000 & 223\\
New cleared norm & $\Z$ & 12 & \textbf{2400} & \textbf{215}\\
New integral-radicand norm & $\Z$ & 12 & \textbf{2912} & \textbf{209}\\
Supplemental identical evaluator & $\Z_{>0}$ & 12 & 70 & \textbf{98}\\
Prior identical polynomial & $\Z_{>0}$ & 12 & 70 & 100\\
\bottomrule
\end{tabular}
\caption{Counts for complete delivered circuits. A fixed scalar multiplication
and a square each cost one multiplication. Fixed numeral leaves and reuse
of an already computed value cost no operations. Domains are kept separate.}
\label{tab:main}
\end{table}

The reduction from $4832$ to $2400$ is $50.33\%$ in total degree; the
operation reduction from $248$ to $215$ is $13.31\%$. The companion uses
six fewer operations at higher degree. Both new twelve-integer points
improve the two inspected twelve-integer baselines in both measures.

There is no established scalar meaning of simplicity. The repository
already contains ten-integer equations, nine-positive equations, and
integer quartics with far fewer variables than the direct construction
but very large fixed coefficients. We do not claim a global minimum or a
new record for number of integer variables. The defensible claims are the
particular equations, their proved transformations, and the measured
improvements against identified parents.

The comparison includes recent incoming archives, not only the older
collection README. The synchronized direct archive already contains a
$522$-natural-variable quartic and a $2088$-integer companion; another
incoming direct construction contains $1690$ integer variables.
The direct comparison includes both benchmarks~\cite{sync,batched}.
An earlier experiment improving an $8619$-integer checkpoint was not
selected as the principal result once these newer inputs were inspected.

\section{A literal twelve-integer equation}\label{sec:equation}
\subsection{Coordinates and fixed coefficients}

All twelve independent variables range over all of $\Z$:
\[
 (j,C,z,h,r,v,w,T,c,n,t,s).
\]
Uppercase $C$ and lowercase $c$ are distinct. Every other symbol below is
a polynomial abbreviation, not an additional existential variable.
Use the complete coherent source-34 six-tuple
\[
 (\beta,K_0,M_C,M_F,\Acone,\Bcone),
 \qquad
 \beta=2^{1064921}-1,\quad
 \Acone=2^{1113},\quad\Bcone=2^{666104}.
\]
Appendix~\ref{app:coeff} gives the bounded arithmetic recipe for all six
numerals. The archive includes the entire finite tile table, machine table,
recipe, and actual integer payloads. $M_F$ already includes its prescribed
shift. None of these numerals is a free program index or a value chosen by
searching for a proof.

\subsection{The arithmetic kernel}
Define the following polynomial registers:
\begin{align}
 J&=6j+1,& q&=\beta J+1,& W_0&=12w+4,\nonumber\\
 X&=qW_0,& Y&=(2v+1)q^3,& E&=XY,\nonumber\\
 a&=E+Y,& H&=4a+3,& \Delta&=a^2+H=(a+2)^2-1.
 \label{eq:outer}
\end{align}
The transport and rank registers are
\begin{align}
 F&=(K_0+W_0)C-z(q-1),\nonumber\\
 R&=(q(q-F)-C)(q^2-1)+(M_C+qM_F)J,\nonumber\\
 k&=R+1+hE,& \eta&=c-Yk,& \zeta&=k-\eta. \label{eq:transport}
\end{align}
Set
\begin{align}
 G&=(\Acone F)(\Bcone C)-(q-\Acone F-\Bcone C)^2,\nonumber\\
 P&=G\eta\zeta-(t^2-G-\eta\zeta+1)^2,\nonumber\\
 \Dg&=(3P-4)s^2+1,& I&=src^2,\nonumber\\
 \Df&=1+\Delta I^2,& Q&=\Delta^2 I^2,& \delta&=Q-1,\nonumber\\
 V&=\Df(c-R+4IT)-c,& \Cc&=QV^2-1,\nonumber\\
 \Dt&=1+XY^2(XY^2+1)k^2,&
 \Dm&=1+\Delta c^2. \label{eq:kernel}
\end{align}
The symbol $P$ in \eqref{eq:kernel} denotes the joint guard, not the final
polynomial $\primary$.

\subsection{The smaller intersective carrier}
Use the classical quintic
\begin{equation}\label{eq:pi}
 \pi(x)=(x^2+3)(x^3-19)=x^5+3x^3-19x^2-57.
\end{equation}
Its intersectivity is a special case of Hyde, Lee, and
Spearman~\cite{hls}. The contribution here is its use in the new
valuation-separated certificate; the quintic itself is not new.
Write
\[
 b_0=X+ac,\qquad u=HT-b_0,\qquad
 \Phi_H(x)=x^5+3H^2x^3-19H^3x^2-57H^5.
\]
Thus $\Phi_H(x)=H^5\pi(x/H)$ as a polynomial identity, with no division
required to evaluate it. In the formal quadratic algebra $d^2=\Dm$,
the remainder $\Phi_H(u+d)=A+Bd$ is computed by
\begin{align}
 Q_0&=u^2+\Dm+3H^2,&
 C_0&=u(u^2+3\Dm)-19H^3,&
 C_1&=3u^2+\Dm,\nonumber\\
 A&=Q_0C_0+2u\Dm C_1,&
 B&=Q_0C_1+2uC_0,\nonumber\\
 S&=IH^5,& Z&=Sn-A. \label{eq:carrier}
\end{align}
The coefficient $B$ is unrelated to any native radix.

\subsection{The ordinary norm polynomial}
Let $\Fn$ be the universal integer polynomial characterized by
\begin{equation}\label{eq:norm}
 \Fn(Z;A_0,\ldots,A_4)
 =
 \prod_{\epsilon\in\{-1,1\}^5}
 \left(Z-\sum_{i=0}^4\epsilon_i\sqrt{A_i}\right).
\end{equation}
Changing the sign of any formal radical permutes the factors, so every
surviving radical exponent is even. This defines an ordinary polynomial,
not a square-root operation allowed in the delivered circuit.
The primary equation is
\begin{equation}\label{eq:primary}
\boxed{\;
 \primary=
 \delta^{16}\Fn\left(
 Z;\,B^2\Dm,\,9S^2\Df,\,81S^2\Dt,\,729S^2\Dg,\,
 \frac{S^2\Cc}{\delta}
 \right)=0 .
 \;}
\end{equation}
The degree of $\Fn$ in each radicand is at most $16$, so the denominator
clears completely over $\Z$. Section~\ref{sec:clearing} gives an explicitly
division-free interpretation. We also prove $\delta\ne0$ on every
integer assignment, so its rational interpretation creates no exceptional
solutions. The lower-operation companion is
\begin{equation}\label{eq:companion}
\boxed{\;
 \companion=
 \Fn\left(
 Z;\,B^2\Dm,\,9S^2\Df,\,81S^2\Dt,\,729S^2\Dg,\,
 S^2\delta\Cc
 \right)=0 .
 \;}
\end{equation}
Both equations use exactly the same twelve coordinates and fixed tuple.

\section{Why the separating weights can be constant}\label{sec:valuation}
Earlier constructions separated radicals by making each successive
coefficient larger than all preceding terms. Such inequalities require
polynomial weights, whose degrees accumulate in the final norm. Here
the existing congruences provide separation at one fixed prime.

\begin{lemma}[Unique least valuation]\label{lem:separation}
Let $r_1,\ldots,r_m$ be nonzero algebraic numbers with $r_i^2\in\Q$.
Suppose that in every subset containing at least two indices, the least
value of $\vp(r_i^2)$ occurs at exactly one index and is even. Then no
nonempty signed subset sum is zero. Consequently, if a signed sum
containing every $r_i$ is rational, all the $r_i$ are rational.
\end{lemma}
\begin{proof}
Fix a subset with at least two terms. Divide its roots by a rational
power of $p$ so that the uniquely least squared valuation becomes zero.
The distinguished square is a unit of $\Z_{(p)}$; the other squares
belong to $p\Z_{(p)}$. Form the product of the distinguished root minus
every signed sum of the other roots. Sign invariance makes this product
an integer polynomial in the squares. Modulo $p$, it is a power of the
distinguished unit square, hence nonzero. None of its factors vanishes.
A nonzero singleton cannot vanish either.

Now work in the multiquadratic splitting field generated by the roots.
Apply any $\Q$-automorphism to the rational sum and subtract. The roots
whose signs change give a vanishing signed subset sum after division by
$2$. The first part forces every root to be fixed by every automorphism,
and hence to belong to the fixed field $\Q$.
\end{proof}
The lemma allows negative rational radicands. It does not assume the
positivity guard before extracting it. Equivalently, one may group roots
by rational square class and use the fact that rational numbers with a
unique least $p$-adic valuation cannot sum to zero.

\begin{lemma}[Intersectivity and the zero scale]\label{lem:intersective}
The polynomial $\pi$ in \eqref{eq:pi} has a root modulo every positive
integer and no root in a real extension of $\Q$ of degree at most two.
\end{lemma}
\begin{proof}
The factor $x^2+3$ has no real root, and $x^3-19$ is irreducible by
Eisenstein at $19$. A degree-three element cannot belong to an extension
of degree one or two.

If $p\equiv1\pmod3$, a primitive cube root of unity $\omega$ in
$\mathbf F_p$ gives $(2\omega+1)^2=-3$, a simple quadratic root.
This includes $p=19$. If $p\equiv2\pmod3$, cubing permutes
$\mathbf F_p^\times$. Since $p\nmid19$, the cubic has a nonzero simple
root. This also handles $p=2$. Simple-root lifting supplies roots at
every power of these primes.
For $p=3$, write $x=1+3t$. The equation $x^3=19$ becomes
\[
 t+3t^2+3t^3=2.
\]
The left side minus two has root $t=2$ and derivative $1$ modulo $3$.
Its successive lifts give roots modulo every power of $3$. The Chinese
remainder theorem finishes. This is the $n=19$ case of the classical
classification~\cite{hls}.
\end{proof}

\begin{lemma}[Unconditional residues]\label{lem:residues}
For every integer assignment,
\begin{gather*}
 q\equiv-1,\quad W_0\equiv1,\quad X\equiv-1,\quad
 a\equiv H\equiv\Delta\equiv0\pmod3,\\
 \Dm\equiv\Df\equiv\Dt\equiv1,\quad
 Q\equiv0,\quad \delta\equiv\Cc\equiv-1,\quad
 u\equiv B\equiv1\pmod3.
\end{gather*}
Also $H\ne0$, $\Delta>1$, $\Dm\ge1$, and $\delta\ne0$.
\end{lemma}
\begin{proof}
The exponent $1064921$ is odd, so $\beta\equiv1\pmod3$.
The residues follow directly. For the odd remainder,
\[
 B\equiv5u^4+10u^2\Dm+\Dm^2\equiv16\equiv1\pmod3.
\]
Thus $B\ne0$ without a monotonicity argument.
The equation $4a+3=0$ has no integer solution. Since $J$ is odd and
nonzero, $|q|\ge\beta-1\ge16$, $|W_0|\ge4$, $|X|\ge64$, and
$|Y|\ge4096$. Hence $|a|\ge63\cdot4096$, giving $\Delta>1$.
Finally $Q=(\Delta I)^2$ is either zero or greater than one.
\end{proof}

\section{Soundness over unrestricted integers}\label{sec:soundness}

The proof first excludes a zero multiplier. It then applies the valuation
lemma before making any positivity assumption about the guard. This order
is needed because all twelve coordinates are unrestricted integers.

\subsection{The zero-multiplier boundary}

\begin{lemma}\label{lem:nonzero-scale}
Neither \eqref{eq:primary} nor \eqref{eq:companion} has an integer zero
with $S=0$. At every zero, $I,s,r,c$ are nonzero and $Q>1$.
\end{lemma}
\begin{proof}
By Lemma~\ref{lem:residues}, $H\ne0$. Thus $S=0$ implies $I=0$,
$Q=0$, and $\delta=-1$. All scaled radicands except $B^2\Dm$ vanish,
and both final polynomials specialize to
\[
 \bigl(A^2-B^2\Dm\bigr)^{16}.
\]
The two factors of the inner expression over $\Q(\sqrt{\Dm})$ are
\[
 A\pm B\sqrt{\Dm}
 =H^5\pi\!\left(\frac{u\pm\sqrt{\Dm}}{H}\right).
\]
Since $\Dm\ge1$, both arguments belong to a real extension of $\Q$ of
degree at most two. Lemma~\ref{lem:intersective} says neither factor can
vanish. Therefore $S\ne0$. Now $I=src^2$ is nonzero, so each of
$s,r,c$ is nonzero; $Q=(\Delta I)^2>1$ follows from $\Delta>1$.
\end{proof}

\subsection{Root extraction without a sign assumption}

Suppose first that $\primary=0$. The unconditional fact $\delta\ne0$
allows the rational interpretation in \eqref{eq:primary}. In a splitting
field, one factor of the norm vanishes. After division by $S\ne0$,
the five radical terms, up to independent signs, are
\begin{equation}\label{eq:normalized-roots}
 \frac{B\sqrt{\Dm}}{S},\qquad
 \sqrt{\Cc/\delta},\qquad
 3\sqrt{\Df},\qquad9\sqrt{\Dt},\qquad27\sqrt{\Dg}.
\end{equation}
Their signed sum is the rational number $n-A/S$. The first four terms
are nonzero by the residues in Lemma~\ref{lem:residues}. Their squared
$3$-adic orders are, respectively,
\begin{equation}\label{eq:orders}
 -2\vt(S),\qquad0,\qquad2,\qquad4.
\end{equation}
Here $\vt(S)=\vt(I)+5\vt(H)\ge5$. If $\Dg\ne0$, the last
squared order is
\[
 6+\vt(\Dg)\ge6;
\]
it need not be even. If $\Dg=0$, omit the zero term. Every subset
with at least two surviving terms has a unique even minimum, because
any such subset contains one of the first four terms. Thus
Lemma~\ref{lem:separation} makes all terms in
\eqref{eq:normalized-roots} rational. This argument also covers
$\Dg<0$; no realness premise was used.

Since $B/S\ne0$, it follows that $\Dm,\Df,\Dt,\Dg$ have rational
square roots. A rational square root of an integer is an integer, so
there are integers $d,f,\tau,g$ with
\begin{equation}\label{eq:extracted-integer-roots}
 d^2=\Dm,\qquad f^2=\Df,\qquad
 \tau^2=\Dt,\qquad g^2=\Dg.
\end{equation}
The auxiliary root is initially rational; Section~\ref{sec:aux-integral}
will prove it integral.

For $\companion=0$, replace the second term of
\eqref{eq:normalized-roots} by $\sqrt{\delta\Cc}$. It again has
squared $3$-adic order zero, so the same proof supplies an integer
$\widetilde y$ with $\widetilde y^2=\delta\Cc$ and all four roots
in \eqref{eq:extracted-integer-roots}.

\subsection{The strict guards}

The extracted square $g^2=\Dg$ proves $\Dg\ge0$. Since $s$ is a
nonzero integer, $P\le0$ would give
\[
 \Dg=(3P-4)s^2+1\le-3.
\]
Hence $P>0$. We record the elementary guard principle to make this
sign inference explicit.

\begin{lemma}[Joint positivity]\label{lem:joint-positive}
For integers $A_1,A_2$, there exists an integer $t$ such that
\[
 A_1A_2-(t^2-A_1-A_2+1)^2>0
\]
if and only if $A_1>0$ and $A_2>0$.
\end{lemma}
\begin{proof}
The inequality implies $A_1A_2>0$. If both factors were negative, then
$t^2-A_1-A_2+1>|A_1|+|A_2|\ge2\sqrt{A_1A_2}$, a contradiction.
Conversely, put $M=A_1+A_2-1\ge1$. An integer square nearest to the
integer $M$ differs from $M$ by an integer of absolute value strictly
less than $\sqrt M$: between consecutive squares $a^2,(a+1)^2$, the
nearest-square error is at most $a$, and at an endpoint it is zero.
Since $M\le A_1A_2$, this choice proves the strict inequality.
\end{proof}

Applying the forward implication to $A_1=G$, $A_2=\eta\zeta$ yields
\begin{equation}\label{eq:strict-guards}
 G>0,\qquad \eta\zeta>0.
\end{equation}
The determinant guard also restores the signs required by the inherited
signed kernel. Set $U=\Acone F$, $W=\Bcone C$, and
$\alpha=q-F-C$. From $UW>(q-U-W)^2$ one obtains $UW>0$. If $q$
had the opposite sign to $U,W$, the square would be at least
$(|U|+|W|)^2\ge4UW$, impossible. Moreover
\[
 |q|>|U|+|W|-\sqrt{UW}\ge\tfrac12(|U|+|W|).
\]
Since both cone scales are at least two, $F,C,\alpha$ have the same
nonzero sign as $q$.

For clarity, $R>0$ already follows at this point. The authentic fixed
tuple has $0<M_C<\beta<M_F$. If $q>0$, then $J>0$ and
\[
 q(q-F)-C=q\alpha+(q-1)C>0.
\]
Both summands in the definition of $R$ are positive. If $q=-q_0<0$,
write $C=-C_0$, $\alpha=-\alpha_0$, with $C_0,\alpha_0>0$.
Then
\[
 q(q-F)-C=q_0(C_0+\alpha_0)+C_0>0,
 \quad J<0,\quad M_C-q_0M_F<0.
\]
Again both summands defining $R$ are positive. Thus no later rank
argument is being used to justify this preliminary sign fact.

\subsection{The rational auxiliary root is integral}\label{sec:aux-integral}

\begin{lemma}[Adjacent-coefficient descent]\label{lem:adjacent-descent}
Let $Q>1$ be an integer and let $V\in\Z$, $y\in\Q$ satisfy
\[
 QV^2-(Q-1)y^2=1.
\]
Then $y\in\Z$.
\end{lemma}
\begin{proof}
Changing signs makes $V,y$ positive. The norm excludes $V=0$.
Write $y=u/b$ in lowest terms with $b>0$. The equation implies
$b^2\mid Q-1$, so $(Q-1)y$ is integral. If $V>1$, use
\begin{align*}
 V'&=(2Q-1)V-2(Q-1)y,\\
 y'&=(2Q-1)y-2QV.
\end{align*}
This transform preserves the norm, and $V'$ is integral. Because
$2Q-1\equiv1\pmod b$, the reduced denominator of $y'$ is still
$b$. The norm gives $y>V$ and $(Q-1)y<QV$, hence
$-V<V'<V$. The norm also excludes $V'=0$. Taking absolute values
therefore gives descent in the positive integer $V$ without changing
$b$. At $V=1$ the equation gives $|y|=1$, so $b=1$.
\end{proof}

For the primary equation, the rational root $y$ extracted from
$\Cc/\delta$ satisfies the lemma's equation because
$\delta=Q-1$ and $\Cc=QV^2-1$. For the companion, take
$y=\widetilde y/\delta$; its square is again $\Cc/\delta$.
Lemma~\ref{lem:nonzero-scale} supplies $Q>1$ in both cases. Thus the
same integer auxiliary Pell witness is recovered from both polynomials.

\subsection{Recovering the main congruence}

Absorb the signs of a vanishing conjugate factor into the integral
roots. The primary equation gives
\begin{equation}\label{eq:root-sum}
 Sn=A+Bd+S(y+3f+9\tau+27g).
\end{equation}
For the companion, replace $y$ on the right by $\widetilde y$.
Consequently $H^5\mid A+Bd$ in both cases. Since
\[
 A+Bd=\Phi_H(HT+d-b_0),
\]
this restores exactly the missing congruence. Indeed, if $p^e\Vert H$
and $v_p(HT+d-b_0)=b<e$, the leading fifth power in
\[
 \Phi_H(z)=z^5+3H^2z^3-19H^3z^2-57H^5
\]
has valuation $5b$, strictly below all three other terms. Divisibility
by $p^{5e}$ is impossible. Doing this for each prime divisor of $H$
proves $H\mid HT+d-b_0$, and therefore
\begin{equation}\label{eq:gamma-restored}
 \gamma=\frac{d-b_0}{H}\in\Z,\qquad d=X+ac+\gamma H.
\end{equation}
This introduces no independent coordinate; $\gamma$ is reconstructed
from an existing zero.

\subsection{The precise inherited signed-kernel interface}

The recovered equations are
\begin{align}
 d^2-\Delta c^2&=1,& d&=X+ac+\gamma H,\nonumber\\
 f^2-\Delta I^2&=1,& I&=src^2\ne0,\nonumber\\
 (\Delta I)^2V^2-((\Delta I)^2-1)y^2&=1,&
 V&=f^2(c-R+4IT)-c,\nonumber\\
 \tau^2-XY^2(XY^2+1)k^2&=1,&
 G&>0,\quad\eta\zeta>0. \label{eq:restored-kernel}
\end{align}
The generalized signed-kernel proof used by the incoming twelve-variable
construction is preserved as
\path{integer12/proofs/INTERPOLATION12.md}, Sections 5--7, and its
independent arithmetic audit as
\path{integer12/proofs/INDEPENDENT_TWELVE_AUDIT.md}, Sections 2--4.
They specify which ordinary rank, orientation, and native-decoding
theorems are inherited. We give the normalized strong-rank argument here
because it prevents an incorrect parity shortcut.

For an integer $A_*>1$, define the Pell polynomials by
\begin{equation}\label{eq:pell-def}
 \chi_{A_*}(m)+\psi_{A_*}(m)\sqrt{A_*^2-1}
   =(A_*+\sqrt{A_*^2-1})^m,\qquad m\ge0.
\end{equation}
The inherited first/main rank bootstrap, applied to the first and last
Pell equations and the strict guards in \eqref{eq:restored-kernel},
gives
\[
 c_*=|c|=\psi_{A_*}(p),\qquad A_*=|a+2|,\qquad p\ge25.
\]
The normalized strong equation in \eqref{eq:restored-kernel} gives
\[
 |I|=\psi_{A_*}(m),\qquad |f|=\chi_{A_*}(m),\qquad m\ge1.
\]
This uses the usual integral Pell classification for
$A_*^2-1$: $A_*+\sqrt{A_*^2-1}$ is the fundamental positive unit
with integer coefficients, as follows by descent on the first coordinate.

Since $c_*^2\mid |I|$, strong divisibility of the Lucas sequence
$\psi_{A_*}$ gives $p\mid m$. Write $m=p\ell$. Expanding the
$\ell$th power of the unit at index $p$ gives
\[
 \frac{\psi_{A_*}(p\ell)}{c_*}
 \equiv\ell\chi_{A_*}(p)^{\ell-1}\pmod{c_*}.
\]
The left side is divisible by $c_*$, and the Pell identity gives
$\gcd(\chi_{A_*}(p),c_*)=1$. Hence $c_*\mid\ell$ and, in
particular, $pc_*\mid m$. These are the strong-rank inputs needed
by the inherited signed auxiliary descent. They do not require $I/c^2$
to be odd.

For the auxiliary interface put $\rho=\Delta I$. The discarded
exposed quotient coordinates are restored by the integer polynomials
\begin{align*}
 o&=f(c-R+4IT),\\
 j_{\rm aux}&=\Delta s^2r^2c^3(c-R)+4src\Df T.
\end{align*}
Direct expansion gives
\[
 V=fo-c=cj_{\rm aux}-R,
 \qquad \rho^2=\Delta(f^2-1).
\]
Also $\rho^2\equiv1-A_*^2\pmod{|f|}$ and
$\rho^2\equiv0\pmod{c_*}$. In the auxiliary Pell classification,
the first coordinate is $|\rho V|$. Its index is odd because
\[
 \chi_{|\rho|}(2j)\equiv(-1)^j\pmod{|\rho|},
 \qquad |\rho|>1,
\]
whereas $\rho\mid\rho V$. This step likewise needs no oddness of
$\rho$.

The inherited signed auxiliary step-down now yields $p=R$ and
$2n_{\rm first}=R+1$; its main-root residue gives
$X=\operatorname{sgn}(c)2^R$. In particular $R$ is odd, so the
fixed residue $X\equiv2\pmod3$ selects $c>0$ and $X=2^R$.
The inherited first-index sign identity gives $k>0$, and
$c/k=Y+\eta/k$, with $0<\eta/k<1$ and $|Y|>1$, then gives
$Y>0$. Finally $q\mid2^R$, $q\equiv1\pmod\beta$, and the
inherited Mersenne-radix sign lemma give $q>0$. The post-power native
decoder and its fixed blank-input calibration supply a finite syntactic
ZFC refutation. The relevant native proofs are retained under
\path{integer12/inherited/root_fusion/provenance/}, including
\texttt{PROOF44.md} and the \texttt{upstream\_ordinary/} calibration
documents. This is a composition of their ordinary theorems, rather than
a claim that finite circuit checks establish the unbounded compiler.

\begin{remark}
An attempted restoration $r_{\rm old}=(\Delta sr-1)/2$ would be
invalid when $\Delta sr$ is even. The normalized rank argument above
is the reason arbitrary nonzero multipliers are admissible. The proof
does not silently import the older odd-multiplier coordinate chart.
\end{remark}

\section{Completeness and the interpolated residue}\label{sec:completeness}

The only new completeness choice is a residue for the shorter quintic.
We reproduce the interpolation mechanism because its quantifier order
is essential: one strong Pell witness must be fixed before the requested
residue class is selected.

The inherited canonical native construction supplies the outer data and
main Pell roots with
\begin{equation}\label{eq:canonical}
 \begin{gathered}
 A_*=a+2>1\text{ even},\qquad R\equiv3\pmod4,\\
 c=\psi_{A_*}(R)>0,\qquad d=\chi_{A_*}(R),\qquad
 \gamma=(d-X-ac)/H\in\Z,\\
 G>0,\qquad\eta\zeta>0,\qquad \tau^2=\Dt.
 \end{gathered}
\end{equation}
In particular $c\equiv R\pmod4$, since the recurrence gives
$\psi_{A_*}(j+2)\equiv-\psi_{A_*}(j)\pmod4$ when $A_*$ is even.
The source-34 tuple's bounded recipe and the exact canonical hypotheses
are recorded in Appendix~\ref{app:coeff} and the inherited coefficient
proof. This is the canonical family used in the incoming interpolation
theorem, not a restriction imposed on arbitrary zeros.

\begin{lemma}[Full residue interpolation]\label{lem:interpolation}
Let $A_*>1$ be even, $\Delta=A_*^2-1$, $R\equiv3\pmod4$, and
$c=\psi_{A_*}(R)>0$. For every prescribed $N>0$, there is a fixed
positive pair $(I,f)$ such that
\[
 f^2-\Delta I^2=1,\qquad4Nc^2\mid I,
\]
and, for this same pair, every residue class $a_0\pmod I$ is represented
by an integer $T$ for which
\begin{equation}\label{eq:interpolation-equations}
 V=f^2(c-R+4IT)-c,\qquad
 (\Delta I)^2V^2-((\Delta I)^2-1)y^2=1
\end{equation}
has an integer solution $y$.
\end{lemma}
\begin{proof}
The Pell recurrence matrix has determinant one, so it has finite order
modulo $4Nc^2$. Choose a positive even multiple $m$ of that order and
put
\[
 I=\psi_{A_*}(m),\qquad f=\chi_{A_*}(m).
\]
Then $4Nc^2\mid I$, $I$ is even, $f$ is odd, and $\gcd(f,I)=1$.

We first justify the valuation fact
\begin{equation}\label{eq:lucas-valuation}
 v_p(I)\ge v_p(m)\qquad(p\mid I).
\end{equation}
For an odd prime dividing $\Delta$, the recurrence reduces to
$\psi_{A_*}(j)\equiv jA_*^{j-1}\pmod p$, so the first zero
index is $p$. For an odd prime not dividing $\Delta$, the norm-one
group in the split quadratic algebra or quadratic field over
$\mathbf F_p$ has order dividing $p-1$ or $p+1$, respectively; the
first zero index is therefore prime to $p$. If $p\mid\psi_{A_*}(j)$,
binomial expansion of the corresponding norm-one unit shows that
multiplying $j$ by $p$ raises the second-coordinate valuation by
exactly one, while multiplying it by a number prime to $p$ leaves
that valuation unchanged. In the first-zero case the initial valuation
is at least one, proving \eqref{eq:lucas-valuation} in both odd-prime
cases. At $p=2$, because $A_*$ is even and $m$ is even, the doubling
identity gives
\[
 v_2(\psi_{A_*}(m))=v_2(m)+v_2(A_*).
\]
This proves the remaining case, and hence $\gcd(m,I^2)\mid I$.

Write $\alpha=A_*+\sqrt\Delta$, so
$\alpha^m=f+I\sqrt\Delta$. In the integer-coefficient quadratic
algebra modulo $f^2$, its norm-one relation gives
\[
 \alpha^{4m}\equiv1-4f\alpha^m\pmod{f^2},
 \qquad \alpha^{4mf}\equiv1\pmod{f^2}.
\]
The second congruence follows by raising the first to the positive
power $f$. Thus $\lambda=4mf$ is a period of both Pell coordinates
modulo $f^2$. Since $\gcd(f,I)=1$ and
\eqref{eq:lucas-valuation} holds,
\[
 \gcd(\lambda,4I^2)\mid4I.
\]
For any specified class $a_0\pmod I$, the two congruences
\begin{equation}\label{eq:interpolation-crt}
 L\equiv R\pmod\lambda,\qquad
 L\equiv R-4Ia_0\pmod{4I^2}
\end{equation}
are therefore compatible. Choose a positive solution $L$. It is
$3\pmod4$. Put $\rho=\Delta I$ and
\[
 V_0=\chi_\rho(L)/\rho,\qquad y=\psi_\rho(L).
\]
The quotient $V_0$ is integral because the index is odd.

For completeness, the exact odd Chebyshev polynomials
$Q_j(X)=\chi_b(2j+1)/b$, $X=b^2$, satisfy
\[
 Q_0=1,\quad Q_1=4X-3,\quad
 Q_{j+1}=(4X-2)Q_j-Q_{j-1}.
\]
Induction gives
\[
 Q_j(X)\equiv(-1)^j(2j+1)\pmod{4X},\qquad
 Q_j(1-A_*^2)=(-1)^j\psi_{A_*}(2j+1).
\]
Since $\rho^2=\Delta(f^2-1)\equiv1-A_*^2\pmod{f^2}$,
these identities and \eqref{eq:interpolation-crt} imply
\begin{equation}\label{eq:two-aux-residues}
 V_0\equiv-c\pmod{f^2},\qquad
 V_0\equiv-R+4Ia_0\pmod{4I^2}.
\end{equation}
Here $L\equiv3\pmod4$ makes $(L-1)/2$ odd; the period $\lambda$
identifies $\psi_{A_*}(L)$ with $c$ modulo $f^2$.

Because $c-R$ is divisible by four and $f^2=1+\Delta I^2$,
the integer
\[
 N_0=V_0+c-f^2(c-R)
\]
is divisible by $f^2$ and congruent to $4Ia_0\pmod{4I^2}$.
The coprimality $\gcd(f^2,4I^2)=1$ now gives
\[
 T=\frac{N_0}{4If^2}\in\Z,\qquad T\equiv a_0\pmod I,
\]
where the last step also uses $f^2\equiv1\pmod I$.
Rearrangement supplies the first equation in
\eqref{eq:interpolation-equations}; the Pell identity for
$\chi_\rho(L),\psi_\rho(L)$ supplies the second. The initial pair
$(I,f)$ depended only on $N$ and the canonical data, not on $a_0$.
\end{proof}

\begin{remark}
The squared-modulus period is $4mf$, not $4m$. The latter is generally
incorrect. Also, even $I$ and odd $f$ are choices made for completeness;
soundness did not restrict $I$ to be even. These two details are recorded
explicitly in the inherited independent arithmetic audit.
\end{remark}

We can now finish completeness for both new equations. Starting with a
finite ZFC refutation, use the inherited canonical data
\eqref{eq:canonical}. Lemma~\ref{lem:joint-positive} selects $t$
with $P>0$. If $P=1$, choose $s=1,g=0$. If $P\ge2$, then
$3P-4>0$ is $2\pmod3$ and hence nonsquare; ordinary Pell existence
gives positive $s,g$ with
\[
 g^2-(3P-4)s^2=1.
\]
Apply Lemma~\ref{lem:interpolation} with $N=s$, then put
$r=I/(sc^2)\in\Z$. At this stage the pair $(I,f)$ and the coordinate
$r$ are fixed.

By intersectivity, choose $a_0$ with
$\pi(a_0+\gamma)\equiv0\pmod I$. The interpolation lemma supplies
$T\equiv a_0\pmod I$ and its integer auxiliary root $y$, without
changing $I,f,r$. The main canonical root $d$ satisfies
\[
 A+Bd=\Phi_H(HT+d-b_0)=H^5\pi(T+\gamma),
\]
so $(A+Bd)/S$ is an integer. Choose
\begin{equation}\label{eq:complete-final-coordinate}
 n=\frac{A+Bd}{S}+y+3f+9\tau+27g.
\end{equation}
One conjugate factor in \eqref{eq:primary} then vanishes. For
\eqref{eq:companion}, replace $y$ in
\eqref{eq:complete-final-coordinate} by $\delta y$. Every choice is
made from already fixed data: native coordinates, $t$, $s$, the pair
$(I,f)$ and $r$, the intersective residue, $(T,y)$, and finally $n$.
There is no circular dependence of the modulus on the requested residue.

Together with Section~\ref{sec:soundness}, this proves that each new
polynomial has an integer zero if and only if there is a finite
syntactic ZFC refutation, with the inherited ordinary calibration
identified above. The elementary arithmetic arguments and those
calibration proofs are formalizable in ZFC, yielding
\eqref{eq:target}. This does not supply a new machine-checked internal
ZFC derivation of the complete biconditional.

\section{The literal denominator-free norm}\label{sec:clearing}

We give a finite ring recipe for the denominator clearing in
\eqref{eq:primary}. This also specifies exactly what the executable
\texttt{norm5.py} computes. Use generic arguments
$x,a_1,a_2,d_0,e_0,N_0,\delta_0$, and work temporarily in the
free rank-four algebra with basis $1,u_0,v_0,u_0v_0$ and relations
$u_0^2=d_0$, $v_0^2=e_0$. For a four-tuple $z=(z_0,z_1,z_2,z_3)$,
its square has coordinates
\begin{align*}
 [z^2]_0&=z_0^2+d_0z_1^2+e_0z_2^2+d_0e_0z_3^2,\\
 [z^2]_1&=2(z_0z_1+e_0z_2z_3),\\
 [z^2]_2&=2(z_0z_2+d_0z_1z_3),\\
 [z^2]_3&=2(z_0z_3+z_1z_2).
\end{align*}
Let $\mathbf e=(1,0,0,0)$ and define
\begin{align*}
 \mathbf t&=(x^2+d_0+e_0,-2x,-2x,2),\\
 \mathbf q&=\delta_0\bigl(\mathbf t-(a_1+a_2)\mathbf e\bigr)
                      -N_0\mathbf e,\\
 r_0&=\delta_0^2a_1a_2,\qquad s_0=\delta_0N_0,\\
 \mathbf h&=\mathbf q^2-4\bigl(r_0+s_0(a_1+a_2)\bigr)\mathbf e,\\
 \mathbf f&=\mathbf h^2-64r_0s_0\mathbf t.
\end{align*}
If $\mathbf f=(f_0,f_1,f_2,f_3)$, put
\begin{align*}
 g_0&=f_0^2+e_0f_2^2-d_0(f_1^2+e_0f_3^2),\\
 g_1&=2(f_0f_2-d_0f_1f_3),\\
 \mathcal C(x;a_1,a_2,d_0,e_0,N_0,\delta_0)&=g_0^2-e_0g_1^2.
\end{align*}
Every operation in this recipe is addition, subtraction, or
multiplication in the coefficient ring. No root, division, or
data-dependent loop occurs in the exported arithmetic circuit.

The three-root identity
\[
 F_8(x;a,b,c)
 =\bigl((x^2-a-b-c)^2-4(ab+ac+bc)\bigr)^2-64abcx^2
\]
shows that $\mathbf f$ represents
$\delta_0^4F_8(x-u_0-v_0;a_1,a_2,N_0/\delta_0)$.
The final two quadratic norms therefore give the universal polynomial
identity
\begin{equation}\label{eq:cleared-norm-identity}
 \mathcal C
 =\delta_0^{16}\Fn(x;a_1,a_2,d_0,e_0,N_0/\delta_0).
\end{equation}
The right side is interpreted after clearing denominators, so the identity
is a polynomial identity even at $\delta_0=0$.
For the principal construction substitute
\[
 (x,a_1,a_2,d_0,e_0,N_0,\delta_0)
 =(Z,B^2\Dm,9S^2\Df,81S^2\Dt,729S^2\Dg,S^2\Cc,\delta).
\]
For the companion, replace the last two arguments by
$(S^2\delta\Cc,1)$.

The independent exact checker constructs $F_{32}$ by successive binomial
translations $F(x+\sqrt a)F(x-\sqrt a)$ and compares every integer
coefficient after denominator clearing. All $20349$ terms match
\eqref{eq:cleared-norm-identity}. Its intermediate term counts are
$2,6,35,495,20349$. This is an exact universal polynomial check;
the additional modular evaluations in the release serve as transcription
controls for the specialized front end.

\section{Attained degree and arithmetic cost}\label{sec:degree}

The degrees of the relevant registers are
\begin{center}\small
\begin{tabular}{@{}lrrrrrrrrrr@{}}\toprule
Register &$S$&$A$&$B$&$\Dm$&$\Df$&$\Dt$&$\Dg$&$\delta$&$V$&$\Cc$\\
\midrule
Degree &34&35&28&14&20&34&46&32&25&82\\\bottomrule
\end{tabular}
\end{center}
The norm $\Fn$ is even in $Z$ and homogeneous of degree $16$ in
$Z^2$ and the five radicands. Consequently its cleared version is a
homogeneous polynomial of degree $16$ in
\[
 \delta Z^2,\quad \delta B^2\Dm,\quad9\delta S^2\Df,
 \quad81\delta S^2\Dt,\quad729\delta S^2\Dg,\quad S^2\Cc.
\]
These six input degrees are, respectively,
\[
 102,\quad102,\quad120,\quad134,\quad146,\quad150.
\]
The coefficient of the sixteenth power of the last input is one.
Therefore $(S^2\Cc)^{16}$ is the entire degree-$2400$ contribution;
every other contribution has smaller degree. This yields the degree
upper bound and reduces attainment to a single leading-form check.

Write
\begin{align*}
 q_0&=6\beta j,& W_{00}&=12w,\\
 a_0&=2q_0^4W_{00}v
     =2^7 3^5\beta^4j^4wv,& I_0&=src^2,\\
 R_0&=-q_0^3(W_{00}C-q_0z),& b&=4I_0T-R_0.
\end{align*}
The leading forms of $\Delta,H,Q,\Df,V$ are respectively
$a_0^2,4a_0,a_0^4I_0^2,a_0^2I_0^2,a_0^2I_0^2b$.
It follows that the leading form of $\primary$ is exactly
\begin{equation}\label{eq:primary-leading}
 2^{320}a_0^{288}I_0^{128}b^{32}.
\end{equation}
Selecting the maximal $T$ power in $b^{32}$ gives the explicit coefficient
\begin{equation}\label{eq:primary-coefficient}
 \bigl[j^{1152}w^{288}v^{288}s^{160}r^{160}c^{320}T^{32}\bigr]
 \primary
 =2^{2400}3^{1440}\beta^{1152}>0.
\end{equation}
The exponents sum to $2400$. Since no other fixed port occurs in this
coefficient, specializing to the actual six-tuple cannot make it vanish.

For the companion, the radicand $S^2\delta\Cc$ has degree $182$,
uniquely above the other radicands and $Z^2$. Its sixteenth power is
the degree-$2912$ part, with leading form
\[
 2^{320}a_0^{352}I_0^{160}b^{32}.
\]
In particular,
\begin{equation}\label{eq:companion-coefficient}
 \bigl[j^{1408}w^{352}v^{352}s^{192}r^{192}c^{384}T^{32}\bigr]
 \companion
 =2^{2848}3^{1760}\beta^{1408}>0.
\end{equation}
The exponents sum to $2912$. These are attained degrees of the
specialized polynomials, not only syntactic degree estimates for circuits
with generic coefficient ports.

The final live-node census is
\begin{center}
\begin{tabular}{@{}lrrrr@{}}\toprule
Polynomial & Additions & Multiplications & Subtractions & Total\\
\midrule
$\primary$ &44&139&32&215\\
$\companion$ &44&133&32&209\\\bottomrule
\end{tabular}
\end{center}
All declared variables and all six coefficient leaves are live. The
counts include the complete denominator clearing and the complete norm,
not just the arithmetic front end. A multiplication by a fixed scalar
is counted, as is each square; already computed registers may be reused.
The exporter removes unused registers and coalesces identical operations.
No assertion of optimal circuit complexity is made.

\section{A direct-calculus quartic in 1190 integers}\label{sec:direct}

\subsection{The literal endpoint and its source}

The second construction improves a different part of the frontier. It
starts from the received synchronized equation in 522 nonnegative
coordinates~\cite{sync}, which directly encodes the numerical calculus
$C_{222}$. It uses neither K215 nor Jones's universal polynomial. All
coefficients of the final polynomial are ordinary small integers.

Let $R_0,\ldots,R_{599}$ be the entire ordered list of quadratic residuals
in \path{direct/output/equation.json}, on the stated coordinates
$x_0,\ldots,x_{1189}$. Define the one ordinary polynomial
\begin{equation}\label{eq:direct1190}
 Q_{\mathrm{direct}}(x_0,\ldots,x_{1189})
   =\sum_{j=0}^{599}R_j(x_0,\ldots,x_{1189})^2,
 \qquad x_i\in\Z.
\end{equation}
The adjacent file \texttt{expanded\_polynomial.json} gives this same
polynomial as a complete list of 26,411 collected monomials. A term
$[c,[i_1,\ldots,i_s]]$ means $c\prod_{r=1}^{s}x_{i_r}$; the empty index
list is the constant monomial. Every coefficient is a literal integer.
There is no remaining power with variable exponent, factorial, binomial,
variable-length product, or unspecified coefficient.

\begin{theorem}[Direct integer endpoint]\label{thm:direct1190}
The polynomial in \eqref{eq:direct1190} has exact degree four, uses all
1190 displayed coordinates, and has coefficient height $83\,900\,628$.
The largest absolute coefficient in any residual is $4556$.
Composing the arithmetic equivalences below with the ordinary inherited
$C_{222}$ calibration proves
\[
 \ZFC\vdash
 \left(\neg\exists\boldsymbol x\in\Z^{1190}\;
                   Q_{\mathrm{direct}}(\boldsymbol x)=0\right)
 \longleftrightarrow\Con(\ZFC).
\]
\end{theorem}

The frozen natural parent has 522 coordinates, 400 quadratic residuals,
20,225 collected monomials, and height $83\,900\,628$. Its exact source
file and the complete received package are included under
\texttt{direct/source/} and \texttt{direct/upstream/synchronized/}.
The new reduction preserves all 522 coordinates and 400 residuals until
the integer-domain conversion is complete. In particular, it modifies
no proof-calculus rule, shared Pell frame, CRT relation, or coefficient
of the natural parent.

The parent represents a finite $C_{222}$ derivation through a commutative
Horn flow, a finite accumulator path, CRT certificates, and synchronized
Pell graphs. Its global proof constructs compatible witnesses for a
common binomial radix, one power frame, and common auxiliary Pell
parameters. Those simultaneous-existence theorems are inherited. Our
new transformation begins after the full natural polynomial has been
fixed; it does not identify independently chosen witnesses.

\subsection{A finite certificate for the necessary signs}

Replacing every one of the 522 natural coordinates by an individually
guarded integer gives the received 2088-integer companion. Many of those
guards are redundant because the equations already imply the required
signs. The following elementary rule makes the redundancy explicit.

\begin{lemma}[Scaled sign inference]\label{lem:directsign}
Suppose a residual is
\[
 c x+Q(\boldsymbol y)=0,\qquad c\in\Z\setminus\{0\},
\]
where $x$ does not occur in $Q$ and every coefficient of $-cQ$ is
nonnegative. If all variables occurring to an odd power in a monomial
of $Q$ are nonnegative, then $x\ge0$.
\end{lemma}
\begin{proof}
The residual implies $c^2x=-cQ$. Every even power of an integer is
nonnegative. The hypotheses on odd powers and coefficients therefore
make the right side nonnegative. Since $c^2>0$, we obtain $x\ge0$.
\end{proof}

This rule proves a sign; it does not eliminate $x$ by division or claim
an integral polynomial formula for $x$. Its monic special case is the
usual implication from $x=F(\boldsymbol y)$ when $F$ has nonnegative
coefficients. The actual certificate has one useful nonmonic instance:
a row equivalent to $16A=20+4T$ proves $A\ge0$ from $T\ge0$. This saves
one further three-coordinate guard beyond the monic-only construction.

The delivered certificate selects 234 source coordinates for explicit
guards, then gives 288 applications of Lemma~\ref{lem:directsign}.
Each record specifies the target, the original residual, the multiplier
$c^2$, the entire polynomial $-cQ$, and its odd-power dependencies.
The records are ordered so that every needed sign is already proved.
In particular, circular implications cannot serve as unconditional
nonnegativity arguments.

An independent enumeration checks that the 234 guarded coordinates
are exactly those that occur as the head of no eligible rule of
Lemma~\ref{lem:directsign} in the literal residual list. Their closure
reaches all 522 coordinates. The selected guard count is therefore
minimal for this specific single-row inference system. That finite
statement says nothing about stronger multi-row arguments, sign
normalizations, or a different presentation of the polynomial.

\begin{lemma}[Quadratic nonnegativity guard]\label{lem:directguard}
For every $x\in\Z$,
\[
 x\ge0\quad\Longleftrightarrow\quad
 \exists a,b,c\in\Z\quad x=a^2+b^2+c(c+1).
\]
\end{lemma}
\begin{proof}
The reverse implication follows from the nonnegativity of squares and
of the product of two consecutive integers. For the forward direction,
$4x+1$ is $1$ or $5$ modulo $8$ and is not divisible by $4$.
Legendre's three-square theorem represents it as a sum of three integer
squares. Modulo $4$, exactly one representing square is odd. Reorder
the terms and write
\[
 4x+1=(2a)^2+(2b)^2+(2c+1)^2.
\]
Subtract one and divide by four.
\end{proof}

For each of the 234 selected coordinates $x_i$, append the quadratic
residual
\begin{equation}\label{eq:directguard}
 x_i-a_i^2-b_i^2-c_i^2-c_i=0,
\end{equation}
using three fresh unrestricted integer variables. The resulting
checkpoint has
\[
 522+3\cdot234=1224\quad\text{coordinates},\qquad
 400+234=634\quad\text{quadratic residuals}.
\]
Its expanded quartic has 23,590 nonzero monomials and coefficient
height $83\,900\,628$.

\begin{proposition}[Natural/integer existence equivalence]
\label{prop:directdomain}
The 1224-coordinate integer checkpoint has a zero if and only if the
frozen 522-coordinate natural parent has a zero.
\end{proposition}
\begin{proof}
At an integer zero, Lemma~\ref{lem:directguard} makes the 234 guarded
source coordinates nonnegative. Apply the 288 ordered records of the
sign certificate: each is justified by Lemma~\ref{lem:directsign} and
by previously proved hypotheses. Every original coordinate is then
natural. The unchanged first 400 residuals give a natural parent zero.

Conversely, start with any natural parent zero and use
Lemma~\ref{lem:directguard} independently for the 234 selected
coordinates. The resulting choices satisfy all new residuals.
Uniqueness of the added guard witnesses is neither claimed nor needed.
\end{proof}

\subsection{Removing signed coordinates without losing a guard}

\begin{lemma}[Integral monic projection]\label{lem:directprojection}
Suppose an integer residual system includes $v-F(\boldsymbol w)=0$
or its negative, with $F\in\Z[\boldsymbol w]$ independent of $v$.
Substitute $F$ for $v$ in every other residual and delete the defining
row and coordinate. Projection is a bijection on the integer zero sets,
with inverse $v=F(\boldsymbol w)$.
\end{lemma}
\begin{proof}
Every source zero has $v=F(\boldsymbol w)$ and therefore projects to a
zero of the substituted system. Conversely, a target zero extends
uniquely by the displayed value of $v$, which is an integer and satisfies
the defining row and every original residual.
\end{proof}

The received exact projection algorithm~\cite{batched} applies this
lemma while requiring every residual to remain quadratic. Affine
restorations preserve that degree bound automatically. A quadratic
restoration is allowed only when its target occurs only as a scalar
multiple, never in a product with another coordinate, in every other
remaining row. The support cap controls expression size and never
truncates a polynomial. All changed residuals are computed exactly.

The integer checkpoint admits 34 delivered substitutions: 23 affine
and 11 quadratic. No unconstrained coordinate is discarded. Negative
restoration coefficients are permitted over $\Z$. Every guard constraint
is preserved either as a substituted residual or through the exact
restoration formula of a deleted guard row. Its consequence is not lost
merely because its original coordinate is eliminated. Apply all inverse
maps in reverse order to recover the complete guarded checkpoint, then
apply Proposition~\ref{prop:directdomain}.

The projected equation therefore has $1224-34=1190$ coordinates and
$634-34=600$ quadratic residuals. Over $\Z$, a sum of squares vanishes
exactly when every residual vanishes. This proves existence equivalence
between \eqref{eq:direct1190} and the frozen natural parent.

\subsection{Logical calibration and the verification boundary}

The parent's numerical-calculus calibration is
\[
 C_{222}(1) \longleftrightarrow\
 \operatorname{BProv}(\mathrm{ZF},\bot)
 \ \longleftrightarrow\
 \operatorname{BProv}(\ZFC,\bot).
\]
The named interfaces are
\path{KEndpoint.c222_iff_zf_refutation} and
\path{ConstructibleChoice.zf_refutation_iff_zfc_refutation}.
Their source copies and pins are included in the inherited provenance.
The second uses the constructible-universe equiconsistency argument.
The calibration is a dependency of the parent, not a newly proved
property of the 1190-variable polynomial.

Negating the composed existence equivalences proves
Theorem~\ref{thm:direct1190}. All new transformations are finite algebra
and ordinary arithmetic formalizable in ZFC, including its internal
natural and integer domains. There is no appeal to an external
enumeration that works only for standard witnesses. There is also no
new complete Lean/Rocq proof term or encoded object-language ZFC
derivation. The result asserts the equivalence and proves neither
consistency nor existence of a zero.

The independent standard-library checker imports neither the producer
nor the optimizer. It authenticates the frozen parent bytes, checks the
whole unchanged source prefix, all literal guards, the complete ordered
sign certificate, all 34 substitutions and coordinate maps, and both
expanded coefficient dictionaries. It verifies that every coordinate
occurs and that degree four is attained. Four deliberately corrupted
artifacts are rejected without relying on hash failure: an inference
coefficient, a guard sign, a restoration, and an expanded coefficient.
Normal and optimized Python runs pass. A second independent projection
implementation also replays the substitutions.

The complete preserved synchronized release additionally passes its own
replay, checking 88 manifest entries and regenerating 22 equation files.
These finite checks establish the literal artifact identities. They do
not replace the unbounded Pell/CRT theorems or the inherited calculus
calibration. The full transitive proof-assistant dependency closure was
not rebuilt.

\subsection{Exact tradeoffs against the incoming direct equations}

\begin{table}[ht]
\centering\small
\begin{tabular}{@{}lrrrr@{}}
\toprule
Construction & Integers & Rows & Terms & Height\\
\midrule
Received synchronized lift & 2088 & 922 & 26791 & 83901150\\
Received batched endpoint & 1690 & 842 & 21381 & 2649643805758\\
New guarded checkpoint & 1224 & 634 & 23590 & 83900628\\
New signed projection & \textbf{1190} & \textbf{600} & 26411 & \textbf{83900628}\\
\bottomrule
\end{tabular}
\caption{Complete direct-calculus equations over unrestricted integers.
All entries have exact degree four. The two received points are taken
from the recent incoming synchronized and batched archives~\cite{sync,batched}.}
\label{tab:direct}
\end{table}

Compared with the received 1690-integer direct equation, the new endpoint
saves 500 variables ($29.59\%$) and 242 residuals ($28.74\%$), while
adding 5030 monomials ($23.53\%$). Its coefficient height is substantially
smaller. Against the synchronized integer lift of the same 522-variable
natural parent, all displayed size measures improve except the unchanged
degree. The 1224-coordinate checkpoint is retained as a
smaller-expression tradeoff, and the natural parent itself is unchanged.

The contribution is a new explicit reduction and combination of the
received constructions. The three-square theorem, sign inference, and
monic projection are standard facts. We claim neither worldwide
priority nor global simplicity. In particular, other repository
quartics use far fewer variables with very different coefficient
payloads. The improvement here is a smaller complete direct-calculus
quartic with literal coefficients and an independently checkable domain
conversion.

\section{What the contribution changes, and what it imports}\label{sec:dependencies}

\subsection{The two distinct calibration chains}
For the twelve-integer equations the proof has the following form:
\begin{align*}
 \exists\boldsymbol x\in\Z^{12}\;\primary(\boldsymbol x)=0
 &\ \Longleftrightarrow\ 
 \text{the fixed source-34 native instance has a witness}\\
 &\ \Longleftrightarrow\ \operatorname{Halt}(\mathrm{K215},\text{blank})\\
 &\ \Longleftrightarrow\ \exists p\;\operatorname{Proof}_{\ZFC}(p,\bot).
\end{align*}
The first arrow is the new arithmetic reduction plus the identified
signed native kernel. The later arrows are the repository's ordinary
compiler and exact-table machine calibration. The companion equation has
the same chain. The complete inherited directory is retained at
\path{integer12/inherited/root_fusion/}; in particular,
\path{provenance/CALIBRATION_STATUS.md} distinguishes the ordinary
mathematical claim from the encoded proof-assistant target.
The inherited proof sources include the signed native arithmetic,
the finite-window tableau decoder, the source-34 coefficient adapter,
and the exact K215 ordinary calibration. Their original paths,
hashes, notices, and historical dependency pins are preserved.

The direct equation instead has the chain
\begin{align*}
 \exists\boldsymbol x\in\Z^{1190}\;
 Q_{\mathrm{direct}}(\boldsymbol x)=0
 &\ \Longleftrightarrow\ 
 \text{the fixed 522-natural-coordinate equation has a witness}\\
 &\ \Longleftrightarrow\ C_{222}(1)\\
 &\ \Longleftrightarrow\ \exists p\;\operatorname{Proof}_{\ZFC}(p,\bot).
\end{align*}
The first arrow is established by the finite sign certificate and
reversible substitutions in Section~\ref{sec:direct}. The natural
equation's arithmetic compiler and numerical-calculus calibration are
imported with their ordinary proofs and source pins. This route makes no
use of a Turing-machine transition table or a Jones universal polynomial.

The displayed implications are used as ordinary mathematical theorems
formalizable in ZFC on its own natural and integer domains. Composing
and negating them proves \eqref{eq:target}. No argument assumes that every
ZFC theorem is true, or that a finite run has searched all proofs.

\subsection{The formalization boundary is material}
The supplied derivations are ordinary mathematical proofs, supplemented
by finite exact certificates. No new complete Lean or Rocq development
was built. In particular, an externally stated theorem about a machine
and a proof predicate is distinct from a derivation of a quoted
biconditional in the project's internal \texttt{SetTheory.Form} calculus.
The latter requires an additional syntactic translation and formal
composition for all internally coded witnesses.

The archived sources themselves record these distinctions. We preserve
them instead of treating compilation, finite simulations, or hashes as
a substitute for an unbounded theorem. If a reviewer declines to accept
an imported ordinary calibration until independently re-establishing it,
the newly checked conclusion is the exact equivalence with that
calibrated parent. The final $\Con(\ZFC)$ interpretation then retains the
explicit calibration dependency. This is the usual dependency structure
of a mathematical reduction, and it identifies the remaining work for a
fully kernel-checked deliverable.

Existing formal work on Diophantine complexity bounds illustrates a
different scale of machine-checked result~\cite{bayerdavid}. It supplies
context, not a theorem automatically validating the present literal
polynomials or their coefficient instantiations.

\section{A small additional improvement: twelve positive variables}
\label{sec:supplement}
The supplemental directory contains a separate optimization of the
repository's ordinary twelve-positive-variable polynomial. Its variables
remain positive, its exact degree remains $70$, and its zero set is
unchanged. Only the evaluation schedule changes from $100$ to $98$
binary operations.

In the parent's unchanged registers set
\[
 d=c^4,\qquad e=1-Q,\qquad v=V+n,\qquad
 E_{\mathrm a}=QV^2-1,\qquad h=ev^2+E_{\mathrm a},\qquad b=eQ.
\]
Then the following is an identity over the polynomial ring:
\begin{equation}\label{eq:polarized}
 (dh+b)^2-dQ(2ev)^2=(dh-b)^2+4dbE_{\mathrm a}.
\end{equation}
Indeed the difference between the two squared terms is $4dhb$ and
$h-E_{\mathrm a}=ev^2$. The right side reuses $E_{\mathrm a}$ and saves
two operations in the complete circuit.

This is a modest evaluator improvement. Norm polarization is already
used elsewhere in the repository. It is not a new representation
theorem, variable reduction, or degree result. Its complete authentic
coefficient tuple is the \emph{compact44} tuple with
$\beta=2^{590069}-1$, preserved separately under
\path{supplemental/ordinary98/native/}.
It must not be mixed with the source-34 tuple in
Section~\ref{sec:equation}.

\section{Assessment and concrete next research questions}
\label{sec:assessment}

The most transferable observation is the prime-valuation separator.
It replaces size comparisons between radicals by congruence information
that the certificate already carries. Once the least squared valuations
are separated, the argument is insensitive to how large the witnesses
are. It also handles an initially negative guard radicand. The use of
the classical quintic is useful because its quadratic remainder has
$B\equiv1\pmod3$; monotonicity is no longer needed to prove $B\ne0$.

The rational auxiliary root is a second independent improvement. Root
extraction first proves it rational, the adjacent-Pell descent proves it
integral, and the final equation carries only the denominator-cleared
norm. This moves an integrality proof into the arithmetic argument and
reduces the polynomial's degree.

The direct result suggests treating sign obligations as part of the
optimization problem. A natural coordinate whose sign follows from
existing equations needs no separate guard. The finite certificate
makes this saving reviewable, and signed monic projection preserves it
without imposing an incorrect positivity condition on eliminated
coordinates.

These observations were not found in the inspected parent constructions
in the forms used here. That is a repository-relative statement, not a
claim of worldwide priority. The classical carrier, norm invariance,
Legendre's theorem, CRT, Pell identities, and monic substitution are
established mathematics. Further novelty assessment should examine
earlier relation-combining and valuation-separation literature, including
the broader Diophantine complexity tradition~\cite{mr}.

Several specific questions remain open in this packet:
\begin{enumerate}
 \item Can a different congruence arrangement separate the same roots
 with a lower-degree carrier or a lower-degree scale $S$? Every proposed
 replacement must retain intersectivity, exclusion of zero-scale roots,
 and full residue interpolation.
 \item Can the auxiliary Pell denominator be shared more efficiently
 with the main congruence certificate? A shorter circuit alone is not a
 proof that its total degree or zero set is preserved.
 \item Can multi-row sign deductions or coordinate sign changes improve
 the direct quartic beyond the $234$ guards needed by the present
 single-row rule family? The restricted optimality certificate leaves
 these stronger methods entirely open.
 \item Can the arithmetic reductions be formalized against the exact
 internal ZFC syntax, including coefficient binding and inherited
 unbounded calibration? This would strengthen the verification status,
 even if no numerical metric improved.
\end{enumerate}
No prospective bound from these questions is included in the result
table.

\section{Reproduction and package map}\label{sec:reproduction}
Python 3.10 or later and its standard library suffice for the arithmetic
checks. A TeX distribution with the packages named in
\path{report/report.tex} builds the article. From the package root:
\begin{verbatim}
python3 verify.py
python3 -O verify.py
cd report
pdflatex -interaction=nonstopmode -halt-on-error report.tex
pdflatex -interaction=nonstopmode -halt-on-error report.tex
\end{verbatim}
The root verifier regenerates the two twelve-integer circuits and checks
their exact output; checks the full generic norm identity by an
independent coefficient calculation; checks the coefficient payload
hashes, circuit structure and arithmetic front ends; rebuilds the direct
equation and checks its full sign and projection certificates; and
checks the identical-polynomial supplemental optimization. The normal
and optimized runs both use explicit exceptions, not assertions that
can disappear under \texttt{-O}.

The release manifest binds the deliverable files. Checks generate
temporary receipts by default and do not rewrite the retained source
certificates. The archived reports record the completed runs. A compact
map is:
\begin{center}
\small
\begin{tabular}{@{}p{.30\linewidth}p{.64\linewidth}@{}}
\toprule
Path & Contents\\
\midrule
\path{report/} & Article sources, bounded coefficient appendix, and PDF.\\
\path{integer12/} & Two new circuits, builders, independent arithmetic
checks, new proof, baseline circuits, and inherited native dependencies.\\
\path{direct/} & Literal 1190-variable quartic, residuals, 1224-variable
checkpoint, sign/projection certificates, independent checker, and
complete synchronized parent.\\
\path{supplemental/ordinary98/} & Identical twelve-positive polynomial,
98-operation evaluator, proof, and its separate coefficient tuple.\\
\path{reviews/} & Independent assistant reviews of the new mathematics
and literal direct artifacts.\\
\path{verification/} & Release receipts and PDF layout review.\\
\path{SOURCE_PINS.json} & Repository commit and received archive hashes.\\
\path{MANIFEST.json} & Final deliverable file hashes.\\
\bottomrule
\end{tabular}
\end{center}

The primary equations are supplied as exact straight-line integer
circuits because writing their fully expanded coefficient lists would
obscure both the definition and the proof. A straight-line circuit with
fixed integer leaves defines a unique ordinary polynomial. The direct
quartic is supplied in both residual and completely expanded form.

\appendix
\section{Complete bounded coefficient recipe}\label{app:coeff}

The following is the full source-34 finite arithmetic recipe~\cite{coeff},
reproduced from the incoming root-fusion report with its licenses retained. The defining payload and source bindings are preserved under \path{integer12/inherited/root_fusion/coefficients/}.
All ranges are bounded constants. The ordered \(943\) tile records
\((x_{i0},x_{i1},x_{i2},x_{i3})\), in west, north, east, south order,
are supplied in
\path{integer12/inherited/root_fusion/coefficients/recipe/wang943_tiles.json}.
Their ordering is part of the definition: selector \(0\) is Start
and selector \(1\) is Forbidden.
The supplied \texttt{k215.tm} contains all \(430\) working actions
from which the source builder reconstructs the whole ordered list.
Its SHA-256 is
\begin{center}\small\texttt{
fdbcef0825886035c3091c8537a69535bc96f9e265f0ae45863543d40ae70090}.
\end{center}
Thus this finite table is delivered data, not an unspecified machine
or a free coefficient parameter.

Put
\[
b=121,\quad L_0=8801,\quad d=bL_0=1064921,\quad
\rho=2^{121},\quad B_0=2^{1064921}.
\]
For \(0\le s<4\), \(0\le j<10\), define
\[
e_s(j)=o_s+2j,\qquad(o_0,o_1,o_2,o_3)=(1020,1000,980,960),
\]
and \(w_j=2^{10+2j}\) for \(0\le j<44\). Let
\[
\begin{split}
\mathcal E={}&\{0,\ldots,942\}
 \cup\{e_s(j):0\le s<4,\ 0\le j<10\}\\
 &\cup\{1063,1064,1095,1096,2195,3296\}.
\end{split}
\]
It has \(989\) elements and maximum \(3296\).
For each tile \(0\le i<943\), set
\[
c_i=1+\sum_{s=0}^3\sum_{j=0}^9
       \operatorname{bit}_j(x_{is}+512)w_{10s+j}
       +\sum_{j=40}^{43}w_j .
\]
Here \(\operatorname{bit}_j(m)=\lfloor m/2^j\rfloor\bmod2\)
is used only to define the finitely many fixed integer coefficients.
It is not an operation on a quantified Diophantine variable.
The other coefficients are
\[
c_{e_s(j)}=w_{10s+j},\quad
(c_{1095},c_{1096},c_{2195},c_{3296})=(w_{40},w_{41},w_{42},w_{43}),
\quad c_{1063}=c_{1064}=0.
\]
Define
\begin{align*}
\mu&=1022+\sum_{j=0}^{43}w_j,\\
Z_{\rm pad}&=(2^{119}-1)\sum_{e=8792}^{8798}\rho^e
                +(2^{27}-1)\rho^{8799},\\
D_C&=\rho^{40}+\rho^{41}+\rho^{1100}+\rho^{2200}
       +\sum_{e\in\mathcal E}c_e(\rho^{4393-e}+\rho^{5495-e}),\\
D_R&=\rho,\\
M_C&=B_0-1-\sum_{\substack{e\in\mathcal E\\e\ne1}}\rho^e
                   -2\rho^{1063},\\
M_{F,\rm native}
 &=4+\mu(\rho^{4393}+\rho^{5495})
     +\sum_{j=0}^9(\rho^{1000+2j}+\rho^{1021+2j})
     +\rho^{2195}+\rho^{3296}+Z_{\rm pad}.
\end{align*}
The six actual fixed ports are precisely
\[
\boxed{\begin{gathered}
\beta=B_0-1,\quad K_0=D_C+B_0D_R,\quad M_C,\quad
M_F=M_{F,\rm native}+B_0-1,\\
A_{\rm cone}=2^{1113},\qquad B_{\rm cone}=2^{666104}.
\end{gathered}}
\]
All displayed exponents are fixed integers. In particular,
\(M_F\) already includes its shift; shifting it again would change
the equation and invalidate the binding.

\begin{table}[ht]\centering
\begin{tabular}{lrr}\toprule
Fixed numeral & Bit length & Minimal unsigned bytes\\\midrule
\(\beta\)&1\,064\,921&133\,116\\
\(K_0\)&1\,065\,043&133\,131\\
\(M_C\)&1\,064\,921&133\,116\\
\(M_F\)&1\,064\,922&133\,116\\
\(A_{\rm cone}\)&1\,114&140\\
\(B_{\rm cone}\)&666\,105&83\,264\\\midrule
Total&&615\,883\\\bottomrule
\end{tabular}
\caption{Exact fixed payload sizes before compression.
The archive stores deterministic gzip encodings of the minimal
unsigned little-endian bytes and binds their raw hashes.}
\end{table}

The source proof verifies the whole coefficient interface, not just
these sizes. It establishes local tile decoding, carry-free geometry,
whole-cell rotation, mask populations, positive canonical slack, and
an explicit rank controller. In particular
\[
\operatorname{pc}(M_C)+\operatorname{pc}(M_{F,\rm native})=d,\quad
M_C\equiv2\pmod4,\quad M_{F,\rm native}\equiv4\pmod8.
\]
The controller uses
\(d=11^2\cdot13\cdot677\), \(\lambda=2341\), and a proved Boolean
subset coverage lemma; it does not assume that \(d\) is a power of
five. Its construction supplies odd canonical exponents, odd cubic
scale, and the canonical parity promises in \eqref{eq:canonical}.
The complete inherited proof and literal builder are included,
with their internal source pins and licenses preserved. See
\path{integer12/inherited/root_fusion/provenance/COEFFICIENT_PROOF.md}
and
\path{integer12/inherited/root_fusion/coefficients/recipe/source34_builder.py}.
The twelve-variable transformation changes none of these six numerals.



\begin{thebibliography}{12}
\bibitem{hls}
Andrea M. Hyde, Paul D. Lee, and Blair K. Spearman.
\emph{Polynomials $(x^3-n)(x^2+3)$ Solvable Modulo Any Integer}.
American Mathematical Monthly \textbf{121}(4) (2014), 355--358.
\href{https://doi.org/10.4169/amer.math.monthly.121.04.355}
{doi:10.4169/amer.math.monthly.121.04.355}.
The $n=19$ case supplies the classical intersective carrier;
an elementary proof is given here.

\bibitem{mr}
Yuri V. Matiyasevich and Julia Robinson.
\emph{Reduction of an arbitrary Diophantine equation to one in 13 unknowns}.
Acta Arithmetica \textbf{27} (1975), 521--553.
\href{https://doi.org/10.4064/aa-27-1-521-553}
{doi:10.4064/aa-27-1-521-553}.

\bibitem{bayerdavid}
Jonas Bayer and Marco David.
\emph{A Formal Proof of Complexity Bounds on Diophantine Equations}.
In ITP 2025, LIPIcs \textbf{352}, 3:1--3:18.
\href{https://doi.org/10.4230/LIPIcs.ITP.2025.3}
{doi:10.4230/LIPIcs.ITP.2025.3}.

\bibitem{repo}
Vladimir Reshetnikov and ProveIt contributors.
\emph{ProveIt}, repository snapshot
\texttt{85ae4070019efd66ea3e58e9a67d287528094e8b}.
\url{https://github.com/VladimirReshetnikov/ProveIt}.
The supplied source pins record all selected incoming archives.

\bibitem{interpolation}
ProveIt incoming archive.
\emph{ConZFC Twelve Integer Interpolation},
\path{docs/incoming/ConZFC_Twelve_Integer_Interpolation_2026-10-24.zip},
at the snapshot in~\cite{repo}.
The ordinary proof \path{proofs/INTERPOLATION12.md} and independent
audit are preserved under \path{integer12/proofs/}.

\bibitem{rootfusion}
ProveIt root-fusion package, inherited by~\cite{interpolation}.
Ordinary native/compiler proofs, normalized signed arithmetic,
coefficient bindings, and calibration status.
Preserved verbatim under \path{integer12/inherited/root_fusion/}.

\bibitem{coeff}
ProveIt source-34 coefficient construction.
\path{integer12/inherited/root_fusion/provenance/COEFFICIENT_PROOF.md},
with authentic native input files and ordinary source copies in the
same dependency tree. The finite recipe is reproduced in
Appendix~\ref{app:coeff}.

\bibitem{sync}
ProveIt incoming archive.
\emph{ConZFC Synchronized Certificates},
\path{docs/incoming/ConZFC_Synchronized_Certificates_2026-10-24.zip},
at the snapshot in~\cite{repo}. Complete received package preserved
under \path{direct/upstream/synchronized/}.

\bibitem{batched}
ProveIt incoming archive.
\emph{ConZFC Batched Certificates},
\path{docs/incoming/con_zfc_batched_certificates_2026-10-23.zip},
at the snapshot in~\cite{repo}. This supplies the 1690-integer comparison
point and the attributed signed monic projection implementation.

\bibitem{ordinary}
ProveIt ordinary-unit-products research directory, at the snapshot
in~\cite{repo}. Parent evaluator and ordinary proof preserved under
\path{supplemental/ordinary98/}, with the independently pinned compact44
native inputs and inherited notices.
\end{thebibliography}


\end{document}
